% answers_short.tex —— 2018 年·天津卷·理科 · 简短答案速查
% 由 tools/make_short_answers.py 自 answers.tex 的 \gksolhead/\gkanswer 提取，
% 并与 2018/2018_数学真题_天津卷理_答案要点.md 逐题比对后生成，勿手改（改源 MD 后重跑）。

\begin{gknotice}
\textbf{使用说明}：本节只给最终答案便于快速核对，完整推导见解析版。答案为编写者独立推导，\textbf{非官方标准答案}。
\end{gknotice}

\gkpart{选择题}
\begin{center}\begin{tabular}{ccccccccc}
\toprule
题号 & 1 & 2 & 3 & 4 & 5 & 6 & 7 & 8 \\
\midrule
答案 & B & C & B & A & D & A & C & A \\
\bottomrule
\end{tabular}\end{center}

\gkpart{填空题}
\begin{itemize}[leftmargin=2.6em]
  \item[\textbf{9.}] $4-\mathrm i$
  \item[\textbf{10.}] $5/2$
  \item[\textbf{11.}] $1/12$
  \item[\textbf{12.}] $1/2$
  \item[\textbf{13.}] $1/4$
  \item[\textbf{14.}] $4<a<8$
\end{itemize}

\gkpart{解答题}
\begin{itemize}[leftmargin=2.6em]
  \item[\textbf{15.}] (1) $B=\pi/3$；(2) $b=\sqrt7$，$\sin(2A-B)=3\sqrt3/14$
  \item[\textbf{16.}] (1) 甲、乙、丙依次抽取 $3,2,2$ 人；(2) $X=0,1,2,3$ 的概率依次为 $1/35,12/35,18/35,4/35$，$E(X)=12/7$；$P(A)=6/7$
  \item[\textbf{17.}] (1) $MN\parallel$ 平面 $CDE$；(2) $1/\sqrt{10}$；(3) $DP=1/\sqrt3$
  \item[\textbf{18.}] (1) $a_n=2^{n-1}$，$b_n=n$；(2) $T_n=2^{n+1}-n-2$；所给恒等式成立
  \item[\textbf{19.}] (1) $x^2/9+y^2/4=1$；(2) $k=1/2$ 或 $11/28$
  \item[\textbf{20.}] (1) $h$ 在 $(-\infty,0)$ 递减，在 $(0,+\infty)$ 递增；(2) 所给恒等式成立；(3) 所述共同切线存在
\end{itemize}
